% Proof inputs: C_F0079 v4; C_F0068 Part A for maximal-Delta rationality.
\section{Liberation dual system criterion}
\label{sec:liberation-dual-system-criterion}

We prove Theorem~\ref{thm:intrinsic-analytic} in this section.

\begin{lemma}\label{lem:liberation-normalization}
If $X\in\M^d$ admits a liberation dual system, then it admits a normalized liberation dual system.
\end{lemma}
\begin{proof}
Let $(L_1,\ldots,L_d)$ be a liberation dual system for $X$.
Define operators $C,\widetilde L_1,\ldots,\widetilde L_d$, each with domain $\M\Omega$, by
\[
 C\xi:=\sum_{k=1}^dL_k\xi-P_\Omega\xi,
 \qquad \widetilde L_i\xi:=L_i\xi-\frac1dC\xi
 \qquad(\xi\in\M\Omega,\ 1\leq i\leq d).
\]
Since $X_j\M\Omega\subseteq\M\Omega$, the liberation relations give, on $\M\Omega$,
\[
 [X_j,C]
 =\sum_{k=1}^d[X_j,L_k]-[X_j,P_\Omega]
 =\sum_{k=1}^d\delta_{kj}[X_k,P_\Omega]-[X_j,P_\Omega]
 =0.
\]
Consequently, for $\xi\in\M\Omega$ and $1\leq i,j\leq d$,
\[
 [X_j,\widetilde L_i]\xi
 =[X_j,L_i]\xi=\delta_{ij}[X_i,P_\Omega]\xi,
 \qquad
 \sum_{i=1}^d\widetilde L_i\xi
 =\sum_{i=1}^dL_i\xi-C\xi
 =P_\Omega\xi.
\]
Since $\mathcal M\Omega\subseteq\operatorname{dom}L_i^*$, we have for $\eta,\xi\in\M\Omega$,
\[
 \langle\eta,\widetilde L_i\xi\rangle
 =\left\langle
 L_i^*\eta-\frac1d\sum_{k=1}^dL_k^*\eta+\frac1dP_\Omega\eta,
 \xi\right\rangle.
\]
By the definition of the adjoint, \(\eta\in\operatorname{dom}\widetilde L_i^*\), and hence \(\mathcal M\Omega\subseteq\operatorname{dom}\widetilde L_i^*\).
Thus \(\widetilde L_1,\ldots,\widetilde L_d\) satisfy all the defining conditions.
\end{proof}

\begin{lemma}[{see \cite[Problem~III.5.18]{Kat95}}]\label{lem:finite-closure}
Let \(T\) be a densely defined closable operator on a Hilbert space \(H\).
If \(T\) has finite rank, then \(\overline T\in\F(H)\), with the same range and rank as \(T\).
\end{lemma}

\begin{lemma}\label{lem:tensor-l2-recovery}
Let $x\in\M$ and suppose that $x-\lambda$ is regular for every $\lambda\in\C$.
The map $\Psi_x$ of Definition~\ref{def:psi} is injective.
If $k\in L^0(\M)\otimes_{\C}L^0(\M)^{\mathrm{op}}$ satisfies
\[
 \Psi_x(k)\in H\otimes_{\C}H,
\]
then $k\in H\otimes_{\C}H$.
\end{lemma}
\begin{proof}
If $T\in\F(H)$ satisfies $[T,x]=0$, then $\operatorname{ran}T$ is a finite-dimensional invariant subspace for $x$.
A nonzero such subspace would contain an eigenvector, contradicting regularity of $x-\lambda$.
Thus $T=0$, so Remark~\ref{rem:cycle-characterization} gives $\Delta(x)=1$.
Lemma~\ref{lem:tensor-injectivity}, applied to the one-element tuple $(x)$, now shows that $\Psi_x$ is injective.

To prove the remaining assertion, put $U:=L^0(\M)$ and define the right multiplication maps
\[
 \begin{aligned}
 R_x:U&\longrightarrow U,& R_x(u)&:=ux,\\
 R_{x^{\mathrm{op}}}:U^{\mathrm{op}}&\longrightarrow U^{\mathrm{op}},&
 R_{x^{\mathrm{op}}}(v^{\mathrm{op}})&:=(xv)^{\mathrm{op}}.
 \end{aligned}
\]
Thus $\Psi_x=R_x\otimes I-I\otimes R_{x^{\mathrm{op}}}$.
Neither right multiplication map has a nonzero finite-dimensional invariant subspace: such a subspace would contain an eigenvector, contradicting regularity of $x-\lambda$.
Since $x\in\M$, both maps preserve $H$, and their restrictions to $H$ have the same property.

For linear maps $A:V\to V$ and $B:W\to W$ on complex vector spaces, $A\otimes I-I\otimes B$ is injective on $V\otimes_{\C}W$ if $B$ has no nonzero finite-dimensional invariant subspace.
Indeed, write a nonzero kernel vector as $\sum_{r=1}^n v_r\otimes w_r$ with $n$ minimal, so both families are linearly independent.
Projecting the second factor modulo $\operatorname{span}\{w_1,\ldots,w_n\}$ shows that this span is $B$-invariant, a contradiction.
The same conclusion holds with $A$ in place of $B$.

The induced linear maps on the quotient spaces are
\[
 \begin{aligned}
 \widetilde R_x:U/H&\longrightarrow U/H,&
 \widetilde R_x(u+H)&:=ux+H,\\
 \widetilde R_{x^{\mathrm{op}}}:U^{\mathrm{op}}/H&\longrightarrow U^{\mathrm{op}}/H,&
 \widetilde R_{x^{\mathrm{op}}}(v^{\mathrm{op}}+H)&:=(xv)^{\mathrm{op}}+H.
 \end{aligned}
\]
Projecting the first factor of $k$ modulo $H$ gives a vector in the kernel of
\[
 \widetilde R_x\otimes I-I\otimes R_{x^{\mathrm{op}}}
 \quad\text{on }(U/H)\otimes_{\C}U^{\mathrm{op}}.
\]
This map is injective by the preceding observation, so exactness of tensoring over $\C$ gives $k\in H\otimes_{\C}U^{\mathrm{op}}$.
Projecting the second factor modulo $H$ now gives a vector in the kernel of
\[
 (R_x|_H)\otimes I-I\otimes\widetilde R_{x^{\mathrm{op}}}
 \quad\text{on }H\otimes_{\C}(U^{\mathrm{op}}/H).
\]
This map is injective because $R_x|_H$ has no nonzero finite-dimensional invariant subspace.
It follows that $k\in H\otimes_{\C}H$.
\end{proof}

For a liberation dual system $(L_1,\ldots,L_d)$, define linear maps
\[
 \Phi_i:\M\otimes_{\C}\M^{\mathrm{op}}
 \longrightarrow\Hom_{\C}(\M\Omega,H),\qquad
 \Phi_i(c\otimes d^{\mathrm{op}}):=cL_id|_{\M\Omega}.
\]
The liberation commutator relations give, for $z\in\M\otimes_{\C}\M^{\mathrm{op}}$ and $s,t\in\M$,
\begin{equation}\label{eq:liberation-tensor-actions}
 \Phi_i(\Psi_j(z))=\delta_{ij}\Theta(\Psi_i(z)),
 \qquad
 \Phi_i((s\otimes t^{\mathrm{op}})z)=s\Phi_i(z)t
 \quad\text{on }\M\Omega.
\end{equation}
Indeed, the first identity on $c\otimes d^{\mathrm{op}}$ is $c[X_j,L_i]d=\delta_{ij}c[X_i,P_\Omega]d$, and the second follows from multiplication in the opposite factor.

\begin{proposition}\label{prop:liberation-delta}
Let $X=(X_1,\ldots,X_d)\in\M^d$, and suppose that $X_i-\lambda$ is regular for every $1\leq i\leq d$ and every $\lambda\in\C$.
If $X$ admits a liberation dual system, then $\Delta(X)=d$.
\end{proposition}
\begin{proof}
Let $T_1,\ldots,T_d\in\F(H)$ satisfy $\sum_i[\rho(X_i),T_i]=0$, and put $k_i:=\Theta^{-1}(T_i)$.
Equation~\eqref{eq:tensor-commutator} gives $\sum_i\Psi_i(k_i)=0$.
By Lemma~\ref{lem:tensor-clearing}, choose regular $s,t\in\M$ such that $m_i:=(s\otimes t^{\mathrm{op}})k_i\in\M\otimes_{\C}\M^{\mathrm{op}}$ for every $i$.
Equation~\eqref{eq:psi-left-linearity} yields $\sum_i\Psi_i(m_i)=0$.
Applying $\Phi_i$ and equation~\eqref{eq:liberation-tensor-actions}, then using density of $\M\Omega$, we obtain
\[
 \Theta(\Psi_i(m_i))=0.
\]
Hence $\Psi_i(m_i)=0$.
The regularity assumption and Lemma~\ref{lem:tensor-l2-recovery} give $m_i=0$.
Cancellation of $s\otimes t^{\mathrm{op}}$ gives $k_i=0$, and therefore $T_i=0$.
Remark~\ref{rem:cycle-characterization} now gives $\Delta(X)=d$.
\end{proof}

\begin{lemma}\label{lem:tensor-images-l2}
Let $X=(X_1,\ldots,X_d)\in\M^d$, and suppose that $X_i-\lambda$ is regular for every $1\leq i\leq d$ and every $\lambda\in\C$.
Let $(L_1,\ldots,L_d)$ be a liberation dual system for $X$.
If $k_1,\ldots,k_d\in L^0(\M)\otimes_{\C}L^0(\M)^{\mathrm{op}}$ satisfy
\[
 h:=\sum_{i=1}^d\Psi_i(k_i)\in\M\otimes_{\C}\M^{\mathrm{op}},
\]
then, for every $i$,
\[
 k_i\in H\otimes_{\C}H,
 \qquad \overline{\Phi_i(h)}=\Theta(\Psi_i(k_i)).
\]
\end{lemma}
\begin{proof}
Fix $i$, write $h=\sum_r c_r\otimes d_r^{\mathrm{op}}$ with $c_r,d_r\in\M$, and put $S_i:=\Phi_i(h)$ on $\M\Omega$.
The adjoint-domain condition in Definition~\ref{def:liberation-dual-system} gives
\[
 \M\Omega\subseteq\dom S_i^*,\qquad
 S_i^*\eta=\sum_r d_r^*L_i^*(c_r^*\eta)
 \quad(\eta\in\M\Omega).
\]
Thus $S_i$ is closable.
By Lemma~\ref{lem:tensor-clearing}, choose regular $s,t\in\M$ such that
\[
 m_j:=(s\otimes t^{\mathrm{op}})k_j
 \in\M\otimes_{\C}\M^{\mathrm{op}}
 \qquad(1\leq j\leq d).
\]
Applying $\Phi_i$ to $\sum_j\Psi_j(m_j)=(s\otimes t^{\mathrm{op}})h$ and using equation~\eqref{eq:liberation-tensor-actions} gives
\begin{equation}\label{eq:liberation-cleared-operator}
 sS_it=\Theta(\Psi_i(m_i))\quad\text{on }\M\Omega.
\end{equation}
The right-hand side has finite rank and $s$ is injective, so the restriction of $S_i$ to $t\M\Omega$ has finite rank.
Since $t$ has dense range, $t\M\Omega$ is dense in $H$.
This restriction is closable, and Lemma~\ref{lem:finite-closure} makes its closure a bounded finite-rank operator on $H$.
The operator $\overline{S_i}$ extends that closure and therefore equals it.
Set $h_i:=\Theta^{-1}(\overline{S_i})\in H\otimes_{\C}H$.
Extending equation~\eqref{eq:liberation-cleared-operator} by boundedness and density, and using the outer bimodule actions of $\Theta$, we obtain
\[
 (s\otimes t^{\mathrm{op}})h_i
 =\Psi_i(m_i)
 =(s\otimes t^{\mathrm{op}})\Psi_i(k_i).
\]
Cancellation of the regular outside factors gives $\Psi_i(k_i)=h_i\in H\otimes_{\C}H$.
Since every $X_i-\lambda$ is regular, Lemma~\ref{lem:tensor-l2-recovery} gives $k_i\in H\otimes_{\C}H$.
\end{proof}

\begin{proof}[Proof of Theorem~\ref{thm:intrinsic-analytic}]
By Proposition~\ref{prop:liberation-delta}, $\Delta(X)=d$.
We prove $(1)\implies(4)\implies(3)\implies(2)\implies(1)$.

$(1)\implies(4)$.
Fix any fraction pair $Q=(s,a,b,t)$ for $u$.
By Theorem~\ref{thm:delta}\textup{(\ref{item:delta-all-fraction-tensors})}, there are tensors $\ell_i\in L^0(\M)\otimes_{\C}L^0(\M)^{\mathrm{op}}$ satisfying equation~\eqref{eq:all-fraction-tensor-criterion}.
Set
\[
 h:=a\otimes t^{\mathrm{op}}-s\otimes b^{\mathrm{op}}
   =\sum_i\Psi_i(\ell_i).
\]
Lemma~\ref{lem:tensor-images-l2} gives $\ell_i\in H\otimes_{\C}H$ and $\overline{\Phi_i(h)}=\Theta(\Psi_i(\ell_i))$.
By definition, $\Phi_i(h)=-\Gamma[Q,L_i]|_{\M\Omega}$.
The common adjoint-domain condition makes $\Gamma[Q,L_i]$ closable by Lemma~\ref{lem:fraction-domains}.
Its closure extends the bounded, everywhere-defined operator $-\overline{\Phi_i(h)}$ and therefore equals it.
With $T_i:=-\Theta(\ell_i)\in\F(H)$, equation~\eqref{eq:tensor-commutator} now gives
\[
 \overline{\Gamma[Q,L_i]}
 =-\Theta(\Psi_i(\ell_i))
 =[\rho(X_i),T_i].
\]
Thus assertion~\textup{(4)} holds.

$(4)\implies(3)$.
Fix any fraction pair $Q=(s,a,b,t)$ for $u$ and choose $T_i\in\F(H)$ such that
$[\rho(X_i),T_i]=\overline{\Gamma[Q,L_i]}$ for every $i$.
By equation~\eqref{eq:liberation-normalization}, on $\M\Omega$ we have
\[
 \begin{aligned}
  \sum_i[\rho(X_i),T_i]
  &=\sum_i\overline{\Gamma[Q,L_i]}\\
  &=s\Bigl(\sum_iL_i\Bigr)b-a\Bigl(\sum_iL_i\Bigr)t\\
  &=sP_\Omega b-aP_\Omega t
   =\Gamma[Q,P_\Omega].
 \end{aligned}
\]
Both sides are bounded, so density extends this identity to $H$, proving assertion~\textup{(3)}.

$(3)\implies(2)$.
Every affiliated operator admits a fraction pair, so this is immediate.

$(2)\implies(1)$.
This is the implication $\textup{(\ref{item:delta-finite-rank})}\implies\textup{(1)}$ of Theorem~\ref{thm:delta}.

For each fixed fraction pair, uniqueness of the tuple $(T_1,\ldots,T_d)$ follows from Theorem~\ref{thm:delta}.
\end{proof}
